\chapter{L'approccio assiomatico}

La definizione frequentista ha un grande pregio: dice che cosa \emph{significa} una probabilità, e da lì abbiamo
ricavato tutte le regole di calcolo. Ha però anche dei limiti. Richiede di ripetere l'esperimento infinite volte, cosa
che nessuno può fare; la frequenza non è una successione fissata, ma cambia da una serie di prove all'altra; e ci sono
eventi che non si possono ripetere affatto («domani pioverà», «questo disco si guasterà entro un anno»).

Nel 1933 il matematico russo Andrej Kolmogorov propose una strada diversa, che è quella usata ancora oggi. Invece di
dire che cosa \emph{è} la probabilità, si elencano poche \textbf{regole} (gli \term{assiomi}) che una probabilità deve
rispettare, scelte in modo che le frequenze le rispettino. Tutto il resto si \textbf{dimostra} a partire da quelle
regole. Le formule sono le stesse del capitolo precedente, ma ora poggiano su basi solide. Prima degli assiomi, però,
bisogna stabilire con precisione \textbf{quali} sottoinsiemi di $\Omega$ hanno diritto a una probabilità.

% ============================================================
\section{Algebre di eventi}
% ============================================================

Le operazioni fra eventi (unione, complementare, intersezione) devono produrre ancora eventi: se sappiamo calcolare la
probabilità di $A$ e di $B$, vogliamo poter calcolare anche quella di «$A$ oppure $B$» e di «non $A$». La famiglia
degli eventi deve quindi essere \textbf{chiusa} rispetto a queste operazioni: applicandole, non se ne deve uscire.

\dfn{Algebra di eventi}{
  Si consideri una famiglia $\mathcal{E}$ di sottoinsiemi di $\Omega$ (non vuota). $\mathcal{E}$ si chiama
  \term{algebra} di sottoinsiemi di $\Omega$, o \term{algebra di eventi}, se gode di due proprietà:
  \begin{enumerate}[label=\alph*)]
    \item è \textbf{chiusa rispetto all'unione}: se $A_1, A_2 \in \mathcal{E}$, allora $A_1 \cup A_2 \in \mathcal{E}$;
    \item è \textbf{chiusa rispetto alla complementazione}: se $A_1 \in \mathcal{E}$, allora
          $\compl{A_1} \in \mathcal{E}$.
  \end{enumerate}
}

Quando lo spazio è infinito (come i pacchetti in coda al router, $\Omega = \N_0$) ci servono anche unioni di
\textbf{infiniti} eventi. Per esempio «il numero di pacchetti è pari» è l'unione degli infiniti eventi elementari
$\{0\}, \{2\}, \{4\}, \ldots$

\begin{richiamo}{l'unione di infiniti insiemi}
Data una successione di insiemi $B_1, B_2, B_3, \ldots$ (una famiglia \emph{numerabile}: uno per ogni naturale), la loro
unione
\[ \bigcup_{n=1}^{\infty} B_n = B_1 \cup B_2 \cup B_3 \cup \cdots \]
è l'insieme degli elementi che stanno in \textbf{almeno uno} dei $B_n$. Il simbolo $\infty$ in alto non è un numero:
dice solo che l'indice non si ferma mai. Per esempio, con $B_n = \{2n\}$:
\[ \bigcup_{n=0}^{\infty} \{2n\} = \{0\} \cup \{2\} \cup \{4\} \cup \cdots = \{0, 2, 4, 6, \ldots\} , \]
l'insieme dei numeri pari. Allo stesso modo $\bigcap_{n=1}^{\infty} B_n$ contiene gli elementi che stanno in
\textbf{tutti} i $B_n$.
\end{richiamo}

\dfn{Sigma-algebra}{
  Un'algebra $\mathcal{E}$ è una \term{$\sigma$-algebra} se, oltre a essere chiusa rispetto alla complementazione e
  all'unione, è chiusa anche rispetto all'unione di una infinità numerabile di suoi elementi:
  \[ B_1, B_2, B_3, \ldots \in \mathcal{E} \quad\Longrightarrow\quad \bigcup_{n=1}^{\infty} B_n \in \mathcal{E} . \]
}

La lettera greca $\sigma$ (\emph{sigma} minuscola) ricorda la «S» di somma: nella tradizione matematica segnala che
l'operazione si può ripetere un'infinità numerabile di volte. Negli spazi finiti la distinzione non conta (le unioni
di infiniti eventi si riducono a unioni di un numero finito di eventi diversi), quindi ogni algebra è anche una
$\sigma$-algebra.

% ============================================================
\section{Le proprietà delle algebre}
% ============================================================

Dalle due sole proprietà della definizione seguono tutte le altre: un'algebra è chiusa anche rispetto
all'intersezione e alla differenza, e contiene sempre $\Omega$ e $\emptyset$.

\teorema[algebra]{Proprietà delle algebre}{
  Sia $\mathcal{E}$ un'algebra di eventi e siano $A, B \in \mathcal{E}$. Allora:
  \begin{enumerate}[label=\alph*)]
    \item $A \cap B \in \mathcal{E}$ (chiusura rispetto all'intersezione);
    \item $A \setminus B \in \mathcal{E}$ (chiusura rispetto alla differenza);
    \item $\Omega \in \mathcal{E}$ e $\emptyset \in \mathcal{E}$.
  \end{enumerate}
}

\begin{dimostrazione}
\passo{Intersezione, con De Morgan} Per la seconda legge di De Morgan $\compl{A \cap B} = \compl{A} \cup \compl{B}$;
facendo il complementare di entrambi i membri:
\[ A \cap B = \compl{\compl{A} \cup \compl{B}} . \]
Ora si costruisce il secondo membro un passo alla volta, restando sempre dentro $\mathcal{E}$: $A$ e $B$ stanno in
$\mathcal{E}$, quindi anche $\compl{A}$ e $\compl{B}$ (chiusura rispetto al complementare); allora anche la loro
unione $\compl{A} \cup \compl{B}$ (chiusura rispetto all'unione); e infine anche il complementare dell'unione, che è
$A \cap B$.
\passo{Differenza} Per il Teorema~\ref{teo:differenza}, $A \setminus B = A \cap \compl{B}$. Siccome $B \in
\mathcal{E}$, anche $\compl{B} \in \mathcal{E}$; e l'intersezione di due elementi di $\mathcal{E}$ sta in
$\mathcal{E}$ per il passo precedente.
\passo{$\Omega$ e $\emptyset$} $\mathcal{E}$ non è vuota, quindi contiene almeno un evento $A$, e con esso
$\compl{A}$. Allora contiene $A \cup \compl{A} = \Omega$ e il suo complementare $\compl{\Omega} = \emptyset$.
\end{dimostrazione}

\teorema[minima]{La più piccola algebra che contiene un evento}{
  Se $A$ è un evento qualsiasi, la più piccola algebra che contiene $A$ è
  \[ \mathcal{E} = \left\{ A,\ \compl{A},\ \Omega,\ \emptyset \right\} . \]
}

\begin{dimostrazione}
\passo{Quei quattro insiemi devono esserci} Un'algebra che contiene $A$ contiene $\compl{A}$, per la chiusura rispetto
al complementare; contiene $A \cup \compl{A} = \Omega$, per la chiusura rispetto all'unione; e contiene
$\emptyset = \compl{\Omega}$, di nuovo per la chiusura rispetto al complementare. Quindi ogni algebra che contiene $A$
contiene almeno i quattro insiemi.
\passo{E bastano} Resta da verificare che i quattro insiemi formano già un'algebra, senza bisogno di aggiungerne
altri. I complementari restano dentro: $A \leftrightarrow \compl{A}$ e $\Omega \leftrightarrow \emptyset$. Le unioni
pure:
\begin{center}
\renewcommand{\arraystretch}{1.25}
\begin{tabular}{c|cccc}
  $\cup$ & $A$ & $\compl{A}$ & $\Omega$ & $\emptyset$ \\ \hline
  $A$ & $A$ & $\Omega$ & $\Omega$ & $A$ \\
  $\compl{A}$ & $\Omega$ & $\compl{A}$ & $\Omega$ & $\compl{A}$ \\
  $\Omega$ & $\Omega$ & $\Omega$ & $\Omega$ & $\Omega$ \\
  $\emptyset$ & $A$ & $\compl{A}$ & $\Omega$ & $\emptyset$
\end{tabular}
\end{center}
Ogni casella della tabella è uno dei quattro insiemi: la famiglia è chiusa rispetto all'unione e al complementare,
cioè è un'algebra.
\end{dimostrazione}

\begin{figure}[H]
\centering
\begin{tikzpicture}[baseline]
  \foreach \x/\lab/\tipo in {0/$A$/1, 3.4/$\compl{A}$/2, 6.8/$\Omega$/3, 10.2/$\emptyset$/4} {
    \begin{scope}[shift={(\x,0)}]
      \ifcase\tipo\or
        \fill[evid] (0,0) ellipse (0.75 and 0.55);
      \or
        \fill[evid, even odd rule] (-1.3,-0.85) rectangle (1.3,0.85) (0,0) ellipse (0.75 and 0.55);
      \or
        \fill[evid] (-1.3,-0.85) rectangle (1.3,0.85);
      \fi
      \draw[mspigray, rounded corners=2pt] (-1.3,-0.85) rectangle (1.3,0.85);
      \draw[mspiA!85!black, line width=0.7pt] (0,0) ellipse (0.75 and 0.55);
      \node[font=\small] at (0,-1.15) {\lab};
    \end{scope}
  }
\end{tikzpicture}
\caption{I quattro elementi della più piccola algebra che contiene $A$ (in giallo l'insieme indicato sotto).}
\end{figure}

\info{Un'algebra è un «livello di informazione»}{
  Si lancia un dado e un amico ci dice soltanto se il risultato è pari o dispari. Con questa informazione possiamo
  decidere se si è verificato $\{2,4,6\}$, $\{1,3,5\}$, $\Omega$ o $\emptyset$, ma non se si è verificato $\{6\}$:
  gli eventi che riusciamo a «vedere» sono proprio i quattro dell'algebra generata da $A = \{2, 4, 6\}$. Chi invece
  vede il dado può decidere su qualsiasi sottoinsieme di $\Omega$: la sua algebra è l'insieme delle parti
  $\mathcal{P}(\Omega)$, con $2^6 = 64$ eventi. Più informazione vuol dire un'algebra più ricca.
}

% ============================================================
\section{Gli assiomi di Kolmogorov}
% ============================================================

Ora possiamo dire che cos'è una probabilità: una funzione che a ogni evento dell'algebra associa un numero fra 0 e 1,
e che rispetta tre regole.

\dfn{Legge di probabilità}{
  Si definisce \term{legge di probabilità} una funzione con dominio $\mathcal{E}$ e codominio $[0, 1]$,
  \[ P : A \in \mathcal{E} \;\longmapsto\; P(A) \in [0, 1] , \]
  che soddisfi gli \term{assiomi di Kolmogorov}:
  \begin{enumerate}[label=\arabic*)]
    \item \textbf{non negatività}: $P(A) \geq 0$ per ogni $A \in \mathcal{E}$, in simboli $\forall A \in \mathcal{E}$;
    \item \textbf{normalizzazione}: $P(\Omega) = 1$;
    \item \textbf{additività}: se $A$ e $B$ sono incompatibili ($A \cap B = \emptyset$), allora
          $P(A \cup B) = P(A) + P(B)$;
    \item[3a)] \textbf{numerabile additività}: se $\{B_n\}_{n \in \N}$ è una collezione numerabile di eventi a due a due
          incompatibili ($B_i \cap B_j = \emptyset$ per $i \neq j$), allora
          \[ P\left( \bigcup_{n=1}^{\infty} B_n \right) = \sum_{n=1}^{\infty} P(B_n) . \]
  \end{enumerate}
}

\dfn{Spazio di probabilità}{
  La terna $(\Omega, \mathcal{E}, P)$, formata dallo spazio campione, da una $\sigma$-algebra di eventi e da una legge
  di probabilità, si chiama \term{spazio di probabilità}.
}

Il simbolo $\forall$ (una «A» capovolta, da \emph{All}) si legge «per ogni»: «$P(A) \geq 0\ \forall A \in
\mathcal{E}$» vuol dire che la probabilità di qualsiasi evento è non negativa. L'assioma 3a parla di una somma di
infiniti termini, che merita un richiamo.

\begin{richiamo}{le somme infinite (serie)}
Che senso ha sommare infiniti numeri $a_1 + a_2 + a_3 + \cdots$? Si sommano i primi $N$ termini, ottenendo le
\term{somme parziali}
\[ S_N = \sum_{n=1}^{N} a_n , \]
e si guarda dove va $S_N$ quando $N$ cresce. Se le somme parziali hanno un limite $S$, si dice che la \term{serie}
converge e si scrive $\sum_{n=1}^{\infty} a_n = S$:
\[ \sum_{n=1}^{\infty} a_n = \lim_{N \to \infty} \sum_{n=1}^{N} a_n . \]
L'esempio più famoso è la \term{serie geometrica} di ragione $\tfrac12$:
\[ \frac12 + \frac14 + \frac18 + \frac1{16} + \cdots = \sum_{n=1}^{\infty} \left(\frac12\right)^n = 1 . \]
Si vede con un disegno: si prende un segmento lungo 1, se ne colora la metà, poi la metà di quel che resta, poi
ancora la metà di quel che resta\ldots{} La parte colorata non supera mai il segmento, e quella che manca si dimezza
a ogni passo: dopo $N$ passi è colorato $S_N = 1 - (1/2)^N$, che tende a 1.
\begin{center}
\begin{tikzpicture}[x=10cm]
  \foreach \a/\b/\lab/\c in {0/0.5/{$\frac12$}/60, 0.5/0.75/{$\frac14$}/48, 0.75/0.875/{$\frac18$}/36,
                             0.875/0.9375/{$\frac1{16}$}/26, 0.9375/0.96875/{}/18, 0.96875/0.984375/{}/12} {
    \fill[mspiric!\c] (\a,0) rectangle (\b,0.5);
    \draw[white, line width=0.8pt] (\b,0) -- (\b,0.5);
    \node[font=\small] at ({(\a+\b)/2},0.25) {\lab};
  }
  \draw[mspiric!80!black, line width=0.6pt] (0,0) rectangle (1,0.5);
  \node[font=\scriptsize, anchor=north] at (0,-0.05) {0};
  \node[font=\scriptsize, anchor=north] at (1,-0.05) {1};
\end{tikzpicture}
\end{center}
Una serie può anche \emph{non} convergere: $1 + 1 + 1 + \cdots$ cresce senza limite. Quando tutti i termini sono
$\geq 0$, come le probabilità, la somma infinita o è un numero o è $+\infty$; l'assioma 3a dice in particolare che,
per eventi incompatibili, è sempre un numero (non più grande di 1).
\end{richiamo}

\begin{esempio}{Aspettare la prima testa}
Si lancia una moneta onesta finché non esce testa. Sia $B_n$ l'evento «la prima testa esce al lancio $n$-esimo»,
cioè $n - 1$ croci seguite da una testa. I lanci sono indipendenti, quindi
\[ P(B_n) = \underbrace{\frac12 \cdot \frac12 \cdots \frac12}_{n-1 \text{ croci}} \cdot \frac12 = \left(\frac12\right)^n . \]
Gli eventi $B_1, B_2, \ldots$ sono incompatibili (la prima testa esce a un solo lancio) e la loro unione è «prima o
poi esce testa». Per la numerabile additività:
\[ P(\text{prima o poi esce testa}) = \sum_{n=1}^{\infty} P(B_n) = \sum_{n=1}^{\infty} \left(\frac12\right)^n = 1 . \]
Senza l'assioma 3a non potremmo fare questo conto: l'evento è l'unione di \emph{infiniti} eventi, e l'additività
semplice (assioma 3) parla solo di due alla volta.
\end{esempio}

\info{Gli assiomi e le frequenze vanno d'accordo}{
  Gli assiomi non sono stati scelti a caso: sono proprio le proprietà che abbiamo trovato per le frequenze nel
  capitolo precedente. La frequenza di un evento è $\geq 0$, quella di $\Omega$ è 1, e se due eventi sono
  incompatibili il numero di prove in cui si verifica uno dei due è la somma dei due conteggi. Per questo la
  probabilità frequentista è una legge di probabilità nel senso di Kolmogorov, e tutto ciò che dimostreremo dagli
  assiomi vale anche per lei.
}

\info{Come si costruisce una probabilità su uno spazio discreto}{
  Su uno spazio discreto basta assegnare a ogni risultato $\omega$ un numero $p(\omega) \geq 0$, con la sola
  condizione che la somma di tutti sia 1. Per un evento qualsiasi si pone poi
  $P(A) = \sum_{\omega \in A} p(\omega)$, la somma dei pesi dei suoi risultati. Questa $P$ rispetta gli assiomi: è
  $\geq 0$, vale 1 su $\Omega$ e, se due eventi non hanno risultati in comune, i loro pesi si sommano. È esattamente
  quello che abbiamo fatto con il dado truccato, assegnando $5/6$ all'1 e $1/30$ alle altre facce.
}

% ============================================================
\section{Le proprietà che seguono dagli assiomi}
% ============================================================

Ritroviamo ora, \textbf{dimostrandole} dagli assiomi, le regole del capitolo precedente. Lo schema è sempre lo stesso:
si scrive l'evento come unione di pezzi \textbf{incompatibili} e si applica l'additività.

\teorema[complementare]{Eventi complementari ed evento impossibile}{
  Per ogni evento $A$: \quad $P(\compl{A}) = 1 - P(A)$. \quad In particolare $P(\emptyset) = 0$.
}

\begin{dimostrazione}
\passo{$A$ e $\compl{A}$ si spartiscono $\Omega$} $\Omega = A \cup \compl{A}$, e i due eventi sono incompatibili.
\passo{Si applicano gli assiomi} Per la normalizzazione e l'additività:
\[ 1 = P(\Omega) = P(A \cup \compl{A}) = P(A) + P(\compl{A}) \quad\Longrightarrow\quad P(\compl{A}) = 1 - P(A) . \]
\passo{L'evento impossibile} Con $A = \Omega$: $P(\emptyset) = P(\compl{\Omega}) = 1 - P(\Omega) = 1 - 1 = 0$.
\end{dimostrazione}

\teorema[sottrazione]{Sottrazione fra eventi}{
  Per due eventi qualsiasi $A$ e $B$:
  \[ P(A \setminus B) = P(A \cap \compl{B}) = P(A) - P(A \cap B) . \]
  È la probabilità che si verifichi $A$ e non $B$ (cioè $A$ meno $B$).
}

\begin{dimostrazione}
\passo{Si taglia $A$ in due pezzi} $A = A \cap \Omega = A \cap (B \cup \compl{B})$ e, per la proprietà distributiva,
\[ A = (A \cap B) \cup (A \cap \compl{B}) . \]
\passo{I pezzi sono incompatibili} Uno sta dentro $B$, l'altro fuori: $(A \cap B) \cap (A \cap \compl{B}) \subseteq
B \cap \compl{B} = \emptyset$.
\passo{Additività} $P(A) = P(A \cap B) + P(A \cap \compl{B})$, e portando $P(A \cap B)$ dall'altra parte si ottiene la
tesi.
\end{dimostrazione}

\teorema[unione]{Unione di eventi non incompatibili}{
  Per due eventi qualsiasi $A$ e $B$:
  \[ P(A \cup B) = P(A) + P(B) - P(A \cap B) . \]
}

\begin{dimostrazione}
\passo{Si scrive l'unione con pezzi incompatibili} Osserviamo preliminarmente che
\[ A \cup B = A \cup (B \cap \compl{A}) : \]
l'unione è fatta da tutto $A$, più la parte di $B$ che sta fuori da $A$. I due eventi a secondo membro sono
incompatibili (uno sta in $A$, l'altro fuori).
\passo{Additività} Quindi $P(A \cup B) = P(A) + P(B \cap \compl{A})$.
\passo{Si usa la sottrazione} $P(B \cap \compl{A})$ è la probabilità che si verifichi $B$ e non $A$ (cioè $B$ meno
$A$): per il teorema precedente, con i ruoli di $A$ e $B$ scambiati, vale $P(B) - P(B \cap A)$. Sostituendo:
\[ P(A \cup B) = P(A) + P(B) - P(A \cap B) . \]
\end{dimostrazione}

\begin{figure}[H]
\centering
\venn{\fill[mspiA!25] \vennA; \begin{scope}\clip \vennA; \fill[evid] \vennB; \end{scope}}
     {$A = (A \cap B) \cup (A \cap \compl{B})$\\ due pezzi incompatibili}
\hspace{1.5cm}
\venn{\fill[mspiA!25] \vennA; \begin{scope}\clip \vennB; \fill[evid, even odd rule] \vennOmega \vennA; \end{scope}}
     {$A \cup B = A \cup (B \cap \compl{A})$\\ tutto $A$ più il giallo}
\caption{Le due scomposizioni usate nelle dimostrazioni: a sinistra $A$ è tagliato da $B$; a destra all'evento $A$
(azzurro) si aggiunge la parte di $B$ che sta fuori da $A$ (giallo).}
\end{figure}

\begin{esercizio}{Indipendenza dei complementari}
Dimostrare che se due eventi $A$ e $B$ sono indipendenti, lo sono anche i loro complementari:
\[ P(A \cap B) = P(A)\,P(B) \quad\Longrightarrow\quad P(\compl{A} \cap \compl{B}) = P(\compl{A})\,P(\compl{B}) . \]
\svolgimento
Per la legge di De Morgan $\compl{A} \cap \compl{B} = \compl{A \cup B}$: «né $A$ né $B$» è il contrario di «almeno
uno dei due». Quindi, usando il complementare, poi l'unione e infine l'indipendenza:
\[
\begin{aligned}
  P(\compl{A} \cap \compl{B}) &= P(\compl{A \cup B}) = 1 - P(A \cup B) = 1 - \left[ P(A) + P(B) - P(A \cap B) \right] \\
  &= 1 - P(A) - P(B) + P(A)\,P(B) .
\end{aligned}
\]
Ora si raccoglie con pazienza. I primi due termini danno $1 - P(A) = P(\compl{A})$; gli ultimi due hanno in comune
$P(B)$:
\[
  P(\compl{A}) - P(B)\left(1 - P(A)\right) = P(\compl{A}) - P(B)\,P(\compl{A})
  = P(\compl{A}) \left(1 - P(B)\right) = P(\compl{A})\,P(\compl{B}) .
\]
Con lo stesso metodo si vede che sono indipendenti anche $A$ e $\compl{B}$: per la sottrazione,
$P(A \cap \compl{B}) = P(A) - P(A \cap B) = P(A) - P(A)\,P(B) = P(A)\,(1 - P(B)) = P(A)\,P(\compl{B})$. Se sapere che
è successo $B$ non dice nulla su $A$, nemmeno sapere che \emph{non} è successo dice qualcosa.
\end{esercizio}

% ============================================================
\section{Probabilità condizionata e teorema di Bayes}
% ============================================================

Nell'approccio assiomatico la probabilità condizionata e l'indipendenza diventano \textbf{definizioni}: si pone
\[ P(A \mid B) = \frac{P(A \cap B)}{P(B)} \quad (P(B) > 0), \qquad
   A, B \text{ indipendenti} \iff P(A \cap B) = P(A)\,P(B) , \]
cioè esattamente le formule che le frequenze ci avevano suggerito. Ne seguono, come prima, la legge della probabilità
composta e quella della probabilità totale. Un fatto in più, che useremo spesso, è che \emph{condizionare non fa
uscire dalla teoria}.

\teorema[condprob]{La probabilità condizionata è una probabilità}{
  Fissato un evento $B$ con $P(B) > 0$, la funzione che a ogni evento $A$ associa $P(A \mid B)$ è una legge di
  probabilità: soddisfa tutti gli assiomi di Kolmogorov.
}

\begin{dimostrazione}
\passo{Non negatività} $P(A \mid B) = P(A \cap B)/P(B)$ è il rapporto fra un numero $\geq 0$ e un numero $> 0$.
\passo{Normalizzazione} $P(\Omega \mid B) = P(\Omega \cap B)/P(B) = P(B)/P(B) = 1$.
\passo{Additività} Se $A_1 \cap A_2 = \emptyset$, i pezzi $A_1 \cap B$ e $A_2 \cap B$ sono ancora incompatibili e, per la
distributiva, $(A_1 \cup A_2) \cap B = (A_1 \cap B) \cup (A_2 \cap B)$. Quindi
\[ P(A_1 \cup A_2 \mid B) = \frac{P(A_1 \cap B) + P(A_2 \cap B)}{P(B)} = P(A_1 \mid B) + P(A_2 \mid B) . \]
Lo stesso ragionamento vale per una famiglia numerabile di eventi incompatibili.
\end{dimostrazione}

Di conseguenza tutte le regole valgono anche «dentro» un condizionamento: per esempio
$P(\compl{A} \mid B) = 1 - P(A \mid B)$.

\subsection{Il teorema di Bayes}

Partiamo dalla legge della probabilità composta, che scrive la stessa intersezione in due modi:
\[ P(A \cap B) = P(B \mid A)\,P(A) = P(A \mid B)\,P(B) . \]
Dividendo per $P(A)$ si ottiene una formula che \textbf{rovescia} il condizionamento.

\teorema[bayes]{Teorema (o legge) di Bayes}{
  Se $P(A) > 0$ e $P(B) > 0$:
  \[ P(B \mid A) = P(A \mid B)\,\frac{P(B)}{P(A)} . \]
  Se poi $\{E_i\}_{i=1}^{M}$ è una partizione di $\Omega$, per ogni $i$:
  \[ P(E_i \mid A) = \frac{P(A \mid E_i)\,P(E_i)}{\sum_{j=1}^{M} P(A \mid E_j)\,P(E_j)} . \]
}

\begin{dimostrazione}
\passo{Prima forma} Dalla probabilità composta $P(B \mid A)\,P(A) = P(A \mid B)\,P(B)$; si dividono i due membri per
$P(A) > 0$.
\passo{Seconda forma} Si applica la prima con $B = E_i$: $P(E_i \mid A) = P(A \mid E_i)\,P(E_i)/P(A)$. Il
denominatore $P(A)$ si scrive con la legge della probabilità totale (Teorema~\ref{teo:totale}) rispetto alla
partizione degli $E_j$: $P(A) = \sum_j P(A \mid E_j)\,P(E_j)$.
\end{dimostrazione}

Il teorema di Bayes serve quando si conosce la probabilità di un \emph{effetto} data la \emph{causa}, $P(A \mid E_i)$,
e si vuole risalire alla probabilità della causa dato l'effetto osservato, $P(E_i \mid A)$. I tre ingredienti hanno
nomi che ritroveremo nella parte di statistica del corso:
\begin{itemize}
  \item $P(E_i)$ è la probabilità \term{a priori} della causa: quanto la si riteneva probabile \emph{prima} di
        osservare $A$;
  \item $P(A \mid E_i)$ è la \term{verosimiglianza}: quanto la causa $E_i$ rende probabile ciò che si è osservato;
  \item $P(E_i \mid A)$ è la probabilità \term{a posteriori}: quanto è probabile la causa \emph{dopo} aver osservato
        $A$.
\end{itemize}

\begin{esempio}{Il canale binario, al contrario}
Riprendiamo il canale binario del capitolo precedente: si trasmette 0 con probabilità $0{,}6$ e 1 con probabilità
$0{,}4$, e ogni bit viene invertito con probabilità $0{,}1$. Avevamo trovato $P(1\ \text{ric.}) = 0{,}42$. Se si
\emph{riceve} un 1, qual è la probabilità che sia stato \emph{trasmesso} un 1? Per Bayes:
\[
  P(1\ \text{tr.} \mid 1\ \text{ric.}) = \frac{P(1\ \text{ric.} \mid 1\ \text{tr.})\,P(1\ \text{tr.})}{P(1\ \text{ric.})}
  = \frac{0{,}9 \cdot 0{,}4}{0{,}42} = \frac{0{,}36}{0{,}42} = \frac{6}{7} \approx 0{,}857 .
\]
Quindi, ricevuto un 1, c'è ancora una probabilità $1/7 \approx 0{,}14$ che sia stato trasmesso uno 0: più del 10\% di
inversione del canale, perché gli 0 trasmessi sono più numerosi degli 1. Decidere quale bit è stato trasmesso, dato
quello ricevuto, è un problema di \emph{decisione} che affronteremo nella parte di statistica.
\end{esempio}

\begin{esempio}{Il test diagnostico}
Una malattia colpisce l'1\% della popolazione. Un test la rileva nel 99\% dei malati, ma dà un falso allarme nel 5\%
dei sani. Una persona fa il test e risulta \textbf{positiva}: qual è la probabilità che sia malata?

Molti rispondono «99\%». Ragioniamo con 10\,000 persone, ognuna rappresentativa:
\begin{center}
\begin{tikzpicture}[x=1cm, y=0.6cm, nodo/.style={draw, rounded corners=2pt, font=\small, align=center}]
  \node[nodo, fill=mspilight] (tot) at (0,0) {10\,000\\ persone};
  \node[nodo, fill=mspiB!15] (m) at (3.6,1.9) {100 malati};
  \node[nodo, fill=mspiA!12] (s) at (3.6,-1.9) {9\,900 sani};
  \draw (tot.east) -- node[above, sloped, font=\scriptsize] {1\%} (m.west);
  \draw (tot.east) -- node[below, sloped, font=\scriptsize] {99\%} (s.west);
  \node[nodo, fill=evid] (mp) at (7.6,2.8) {99 positivi};
  \node[nodo] (mn) at (7.6,1.0) {1 negativo};
  \node[nodo, fill=evid] (sp) at (7.6,-1.0) {495 positivi};
  \node[nodo] (sn) at (7.6,-2.8) {9\,405 negativi};
  \draw (m.east) -- node[above, sloped, font=\scriptsize] {99\%} (mp.west);
  \draw (m.east) -- node[below, sloped, font=\scriptsize] {1\%} (mn.west);
  \draw (s.east) -- node[above, sloped, font=\scriptsize] {5\%} (sp.west);
  \draw (s.east) -- node[below, sloped, font=\scriptsize] {95\%} (sn.west);
\end{tikzpicture}
\end{center}
I positivi sono $99 + 495 = 594$, e di questi solo 99 sono malati: $99/594 = 1/6 \approx 0{,}17$. Con il teorema di
Bayes, rispetto alla partizione {malato, sano}:
\[
  P(M \mid +) = \frac{P(+ \mid M)\,P(M)}{P(+ \mid M)\,P(M) + P(+ \mid S)\,P(S)}
  = \frac{0{,}99 \cdot 0{,}01}{0{,}99 \cdot 0{,}01 + 0{,}05 \cdot 0{,}99} = \frac{0{,}0099}{0{,}0594} \approx 0{,}167 .
\]
Il test è ottimo, eppure un positivo è malato solo una volta su sei: i sani sono così tanti che i loro falsi allarmi
superano di gran lunga i veri positivi. È l'effetto della probabilità a priori, $P(M) = 0{,}01$, che è piccola.
\end{esempio}

\warning{Non scambiare $P(A \mid B)$ con $P(B \mid A)$}{
  $P(+ \mid M) = 0{,}99$ e $P(M \mid +) \approx 0{,}17$ sono due numeri molto diversi. Confonderli è uno degli errori
  più comuni (e più pericolosi) del ragionamento probabilistico; il teorema di Bayes dice come passare dall'uno
  all'altro, e mostra che serve anche la probabilità a priori.
}

% ============================================================
\section{Perché una probabilità sta fra 0 e 1}
% ============================================================

Nella definizione di legge di probabilità abbiamo scritto che il codominio è $[0, 1]$. In realtà non era necessario
dirlo: lo si può dimostrare dagli assiomi. Il primo passo è un risultato utile anche da solo.

\teorema[monotonia]{Monotonia della probabilità}{
  Se l'evento $A$ è contenuto nell'evento $B$ ($A \subseteq B$), allora $P(A) \leq P(B)$.
}

\begin{dimostrazione}
\passo{Si scrive $B$ con pezzi incompatibili} Se $A \subseteq B$ possiamo scrivere
\[ B = A \cup (\compl{A} \cap B) , \]
cioè $B$ è fatto da $A$ più la parte di $B$ che sta fuori da $A$. I due eventi non si intersecano:
$A \cap (\compl{A} \cap B) = \emptyset$.
\passo{Additività} Per il terzo assioma $P(B) = P(A) + P(\compl{A} \cap B)$.
\passo{Non negatività} Per il primo assioma $P(\compl{A} \cap B) \geq 0$, quindi $P(B) \geq P(A)$.
\end{dimostrazione}

\teorema[zerouno]{La probabilità sta fra 0 e 1}{
  Per ogni evento $A$: \quad $0 \leq P(A) \leq 1$.
}

\begin{dimostrazione}
\passo{} Ogni evento sta fra l'evento impossibile e quello certo: $\emptyset \subseteq A \subseteq \Omega$. Per la
monotonia, e ricordando $P(\emptyset) = 0$ e $P(\Omega) = 1$:
\[ 0 = P(\emptyset) \leq P(A) \leq P(\Omega) = 1 . \]
\end{dimostrazione}

La monotonia ha anche un significato intuitivo: se ogni volta che succede $A$ succede anche $B$, allora $B$ succede
\emph{almeno} altrettanto spesso. Per esempio, «esce un 6» implica «esce un numero pari», e infatti
$P(\{6\}) = 1/6 \leq P(\{2,4,6\}) = 1/2$.

\gbox{In sintesi}{mspiA!80!black}{
\begin{itemize}
  \item Gli eventi formano un'\textbf{algebra}: una famiglia chiusa rispetto all'unione e al complementare (e quindi
        anche all'intersezione e alla differenza, e contenente $\Omega$ e $\emptyset$). Una \textbf{$\sigma$-algebra}
        è chiusa anche rispetto alle unioni numerabili. La più piccola algebra che contiene $A$ è
        $\{A, \compl{A}, \Omega, \emptyset\}$.
  \item \textbf{Assiomi di Kolmogorov}: $P(A) \geq 0$; $P(\Omega) = 1$; additività (anche numerabile) per eventi
        incompatibili. $(\Omega, \mathcal{E}, P)$ è uno \textbf{spazio di probabilità}.
  \item Dagli assiomi: $P(\compl{A}) = 1 - P(A)$, $P(\emptyset) = 0$, $P(A \setminus B) = P(A) - P(A \cap B)$,
        $P(A \cup B) = P(A) + P(B) - P(A \cap B)$, monotonia e $0 \leq P(A) \leq 1$.
  \item Se $A$ e $B$ sono indipendenti, lo sono anche i complementari. La probabilità condizionata a $B$ è a sua
        volta una probabilità.
  \item \textbf{Bayes}: $P(B \mid A) = P(A \mid B)\,P(B)/P(A)$; con una partizione,
        $P(E_i \mid A) = P(A \mid E_i)P(E_i) / \sum_j P(A \mid E_j)P(E_j)$. A priori, verosimiglianza, a posteriori.
\end{itemize}
}
